%======================================================================================================
%
% BEAMER PACKAGES FILE (A. FIORUCCI)
% AFFILIATION : ECOLE POLYTECHNIQUE
%
%======================================================================================================

%======================================================================================================
% LATEX PACKAGES
%======================================================================================================

% Typesetting packages
	\usefonttheme{professionalfonts}
	\usepackage[UKenglish]{babel}
	\usepackage[T1]{fontenc}
	\usepackage{libertine}
	\usepackage[libertine]{newtxmath}
	\usepackage[cal=euler,scr=boondoxo]{mathalfa}
		
%% Math and symbols
	\usepackage{amssymb}
	\usepackage{amsmath}
	\usepackage{amsfonts}
	\usepackage{mathtools}
	\usepackage{bm}
	\usepackage{upgreek}
	\usepackage{cancel}
	\usepackage{fontawesome}
	\usepackage{wasysym}
	\usepackage{bigints}
	\usepackage{slashed}
	
% Layout
	\usepackage[absolute,overlay]{textpos}
	\usepackage{graphicx}
	\usepackage{float}
	\usepackage{graphbox}
	
% Tikz drawings
	\usepackage{xcolor}
	\usepackage{tikz}
	\usetikzlibrary{calc}
	\usetikzlibrary{decorations.pathreplacing}
	\usetikzlibrary{decorations.pathmorphing}
	\usetikzlibrary{decorations.markings}
	\usetikzlibrary{arrows.meta}
	\usetikzlibrary{positioning}
	\usetikzlibrary{shadows,shadings}
	\usepackage{tikzsymbols}
	\usepackage{pgf-pie}
	
% Arrays
	\usepackage{array}
	\newcolumntype{C}[1]{>{\centering\let\newline\\\arraybackslash\hspace{0pt}}m{#1}}	
	
% Page format
	\usepackage[toc,page]{appendix}
	\usepackage{caption}
	\usepackage{changepage}
	\usepackage{setspace}
	\usepackage[labelformat=simple]{subfig}
	\renewcommand{\thesubfigure}{\relax}
	\usepackage[export]{adjustbox}
	\usepackage{appendixnumberbeamer}
	\usepackage{multirow}

% Multimedia content
	\usepackage{multimedia}
	
% References
	\usepackage{hyperref}

	\usepackage{comment}
	
%======================================================================================================
% BEAMER SPECIFIC SETTINGS
%======================================================================================================

% Define institute colors
	\definecolor{xblue}{RGB}{1,66,106}
	\definecolor{xlightblue}{RGB}{184,211,225}
	\definecolor{xgray}{RGB}{150,150,150}
	\definecolor{xsepia}{RGB}{166,139,78}
	\definecolor{xred}{RGB}{200,32,33} %{169,32,33}
	\definecolor{xorange}{RGB}{254,80,0}
	\definecolor{xgreen}{RGB}{67,176,42}
	\definecolor{xpurple}{RGB}{77,22,84}
	\definecolor{xpink}{RGB}{227,28,121}

% Main theme
	\usetheme{Madrid}
	\setbeamersize{text margin left=20pt,text margin right=20pt} 

% Main colour
	\usecolortheme[named=xblue]{structure}

% Layout of the slide
	\setbeamertemplate{navigation symbols}{} % No navigation symbol at the bottom of the page
	\setbeamercolor{frametitle}{fg=xblue,bg=gray!20}
	\setbeamerfont{frametitle}{family=\rmfamily,shape=\scshape}
	
% Layout of the blocks
\setbeamercolor{block title}{fg=xblue,bg=gray!20}
	\setbeamerfont{block title}{family=\rmfamily,shape=\scshape}


% Item bullets
	\setbeamertemplate{enumerate items}[square]
	\setbeamertemplate{sections/subsections in toc}[square]
	
\newcommand{\sectionframe}{
\begingroup
\setbeamercolor{page number in head/foot}{fg=xblue}
\begin{frame}{Outline of the talk}
	\tableofcontents[currentsection]
\end{frame}
\endgroup
\addtocounter{framenumber}{-1}
}

%======================================================================================================
% PERSONAL COMMANDS
%======================================================================================================

\usepackage{appendixnumberbeamer}
\newcommand{\backupbegin}{
   \newcounter{framenumberappendix}
   \setcounter{framenumberappendix}{\value{framenumber}}
}
\newcommand{\backupend}{
   \addtocounter{framenumberappendix}{-\value{framenumber}}
   \addtocounter{framenumber}{\value{framenumberappendix}} 
}

% Custom colours
	\definecolor{darkred}{RGB}{192,0,0} 
	\definecolor{Cgreen}{RGB}{0,153,0}
	\definecolor{oldbook}{RGB}{255,222,117}

% Custom symbols
	\newcommand{\D}{\text{d}}
	\newcommand{\indep}{\raisebox{0.05em}{\rotatebox[origin=c]{90}{$\models$}}}
	\newcommand{\ndelta}{\delta\hspace{-0.50em}\slash\hspace{-0.05em} }

% Semi-direct sum
	\usepackage{pict2e}
	\makeatletter
		\newcommand{\loplus}{\mathbin{\mathpalette\dog@lsemi{+}}}
		\newcommand{\dog@lsemi}[2]{\dog@semi{#1}{#2}{270,90}}
		\newcommand{\dog@semi}[3]{%
  		\begingroup
  			\sbox\z@{$\m@th#1#2$}%
  			\setlength{\unitlength}{\dimexpr\ht\z@+\dp\z@\relax}%
  			\makebox[\wd\z@]{\raisebox{-\dp\z@}{%
    		\begin{picture}(1,1)
    			\linethickness{\variable@rule{#1}}
    			\roundcap
    			\put(0.5,0.5){\makebox(0,0){\raisebox{\dp\z@}{$\m@th#1#2$}}}
    			\put(0.5,0.5){\arc[#3]{0.5}}
    		\end{picture}%
  		}}%
  		\endgroup
		}
		\newcommand{\variable@rule}[1]{%
 		\fontdimen8  
  		\ifx#1\displaystyle\textfont3\else
    	\ifx#1\textstyle\textfont3\else
      	\ifx#1\scriptstyle\scriptfont3\else
        \scriptscriptfont3\relax
  		\fi\fi\fi
		}
	\makeatother

% Text article references
	\newcommand{\cit}[1]{{\footnotesize \color{xsepia} \normalfont{[#1]}}}
% Text talks
	\newcommand{\cittalk}[1]{{\footnotesize \color{xorange} \normalfont{[#1]}}}
	
% Small quotations
	\newcommand{\quot}[1]{{\footnotesize \color{gray} \normalfont{#1}}}
		
% Important results
	\usepackage{mdframed}
	\global\mdfdefinestyle{HL}{
		linecolor=xblue, linewidth=1pt,%
		topline=false, bottomline=false, rightline=false,%
		align=left, backgroundcolor=xblue!20, userdefinedwidth=\textwidth,%
		innerleftmargin=5, innertopmargin=5, innerbottommargin=5, innerrightmargin=5, skipabove=5}
		
	\global\mdfdefinestyle{HLrelax}{
		linecolor=xblue, linewidth=1pt,%
		topline=false, bottomline=false, rightline=false,%
		align=left, backgroundcolor=xblue!20,%
		innerleftmargin=5, innertopmargin=5, innerbottommargin=5, innerrightmargin=5, skipabove=5}

	\newcommand{\hl}[1]{
		\begin{mdframed}[style=HL]
		#1
		\end{mdframed}
		\vspace{8pt}
	}
	
	\newcommand{\hlrelax}[1]{
		\begin{mdframed}[style=HLrelax]
		#1
		\end{mdframed}
		\vspace{8pt}
	}

	
%======================================================================================================